% Template: Tam giác ABC cơ bản
% Dùng cho: hình học phẳng, minh hoạ bài toán tam giác
% --------------------------------------------------------
\documentclass[border=5pt]{standalone}
\usepackage[utf8]{inputenc}
\usepackage[T5]{fontenc}
\usepackage[vietnamese]{babel}
\usepackage{tikz}
\usepackage{helvet}
\renewcommand{\familydefault}{\sfdefault}
% ── Hệ màu tập trung (chỉnh tại đây để đổi theme) ──────────
\usepackage{xcolor}
\definecolor{mauchinh} {HTML}{1565C0}
\definecolor{mauphu}   {HTML}{43A047}
\definecolor{maunoibat}{HTML}{E53935}
\definecolor{mauvien}  {HTML}{424242}
{mauchinh} {HTML}{1565C0}
{mauphu}   {HTML}{43A047}
{maunoibat}{HTML}{E53935}
{mauvien}  {HTML}{424242}
\usetikzlibrary{calc, angles, quotes, arrows.meta}

\begin{document}
\begin{tikzpicture}[
    scale=1.2,
    >=Stealth,
    line join=round,
    line cap=round,
    every node/.style={font=\sffamily\small},
    thick,
]
    % Đỉnh A (trên)
    \path (1, 3) coordinate (A);
    % Đỉnh B (trái)
    \path (0, 0) coordinate (B);
    % Đỉnh C (phải)
    \path (4, 0) coordinate (C);

    % Ba cạnh AB, BC, CA
    \draw (A) -- (B) -- (C) -- cycle;

    % Nhãn đỉnh A
    \path (A) node[above] {$A$};
    % Nhãn đỉnh B
    \path (B) node[below left] {$B$};
    % Nhãn đỉnh C
    \path (C) node[below right] {$C$};

    % Điểm tại các đỉnh
    \foreach \p in {A, B, C} {
        \fill (\p) circle (2pt);
    }
\end{tikzpicture}
\end{document}
